\begin{table}[htbp]
\centering
\caption{\new{IV-DML estimates with and without controls for opaque-channel exposure, dual outcome.}}
\label{tab:t1_dual}
\begin{threeparttable}
\small
\begin{tabular}{lcc}
\toprule
Specification & Gov. outcome & Centrão outcome \\
\midrule
Leg 55 Base spec & $+1.89^{**}$ & $+0.0011$ \\
Leg 55 Base + proxies & $+2.02^{***}$ & $+0.0013$ \\
Leg 56 Base spec & $-0.94^{***}$ & $-0.0006^{**}$ \\
Leg 56 Base + proxies & $-0.94^{***}$ & $-0.0004$ \\
\bottomrule
\end{tabular}
\begin{tablenotes}
\small
\item Notes: Coefficients in percentage points per R\$1M of pre-vote committed amendment (Centrão column reported with four decimals because magnitudes are an order smaller). Base spec is the preferred PLIV-DML with full-clean controls plus party fixed effects, after explicit exclusion of the multi-RP variables and the secret-budget proxies from the control set. Augmented spec adds $d_{\mathrm{rp9}}$ and $\mathrm{share}_{\mathrm{pix}}$ as additional controls. The Bolsonaro-era coefficient with the government outcome is unchanged by the augmentation, indicating that the visible-amendment backfire is identified independently of direct controls for opaque-channel exposure. The Centrão outcome estimates are an order of magnitude smaller; under either specification, visible amendments produce essentially no movement in alignment with the Centrão bloc.
\end{tablenotes}
\end{threeparttable}
\end{table}